\documentclass[11pt]{article}
\usepackage[a4paper,margin=2.6cm]{geometry}
\usepackage[T1]{fontenc}
\usepackage[utf8]{inputenc}
\usepackage{lmodern}
\usepackage{graphicx}
\usepackage{caption}
\usepackage{booktabs}
\usepackage{xcolor}
\usepackage{enumitem}
\usepackage[hidelinks]{hyperref}
\usepackage{parskip}
\graphicspath{{../../current/output/figures/}}
\captionsetup{font=small,labelfont=bf,width=.94\textwidth}
\newcommand{\todo}[1]{\textcolor{red!70!black}{[#1]}}
\newcommand{\fig}[3]{\begin{figure}[htbp]\centering\includegraphics[width=#3\textwidth]{#1}\caption{#2}\end{figure}}
\setlist{nosep}

\title{\textbf{I Can't Tell You Who to Vote For, But\ldots}\\[4pt]
\large What happens when ten AI models meet the Brazilian election}
\author{Letícia Auriemo \and Luca Moreno Louzada}
\date{Draft, 12 September 2026}

\begin{document}
\maketitle

\noindent\small\textbf{Acknowledgements.} We thank the team at Ranqia for partnering with us on the ChatGPT web data collection, and Leandro Consentino for helpful comments.

\noindent\small\textbf{Use of AI tools.} Claude Code (Anthropic) and Codex (OpenAI) were used as coding assistants for writing API call scripts, data cleaning scripts, and preparing tables and figures. The authors designed the experiment, decided on hypotheses to test and results to show, and wrote the piece; AI was used only for language editing.
\normalsize

\vspace{1em}
``Who should I vote for?'' is an ordinary question to ask someone whose judgment you trust. Ask a chatbot and you get an unusually accommodating adviser. You can tell it precisely what matters to you, disagree with its answer, and ask it to adapt.

People already bring political questions to these tools. Andy Hall and Dan Thompson looked at how Americans used Claude around state primary elections and found a distinct uptick in political conversations. While most users asked for explanations or policy breakdowns, about 3\% of conversations ended with the model offering direct advice or recommendations \todo{CITE Free Systems piece}.

Other experiments show what happens when users ask models who to vote for. In Japan, Sho Miyazaki and Andy Hall found that models repeatedly recommended the Japanese Communist Party to hypothetical voters with left-leaning views, even when other parties shared those same policy positions \todo{CITE Japan piece}. A general political stance turned into a narrow recommendation.

Brazil provides the next case study. AI assistants have acquired a large audience just as the country prepares for a national election. Voters must choose among many parties and candidates, leaving an assistant with several distinct ways to translate the same political priority into a recommendation. And unlike most countries, Brazil has adopted a rule that directly addresses what these systems may tell voters.

\section*{Brazil wrote a rule for this}

In March, Brazil's Superior Electoral Court, known as the TSE, prohibited providers of AI systems from ranking, recommending, suggesting or prioritizing candidates and parties, even when a user expressly asks. The provision also covers direct and indirect electoral favoritism through automated answers \todo{CITE Resolução TSE 23.755/2026, art.~28, §1º-C}. Developers responded as the campaign approached: OpenAI, for instance, published a Brazil-specific compliance plan on August 16 and opened a reporting channel for election-related outputs.

Complying with that rule, however, is rarely straightforward. An assistant can state that it cannot tell someone how to vote, then immediately identify the candidate who ``most closely matches'' the user's profile. It can offer a personalized shortlist or recommend a party instead of an individual. Without ever saying ``vote for X'', the model still narrows the voter's options.

This gray area faces a real-world test in Brazil's October elections, when voters select a president, congressional representatives, and state governors. In the presidential race, if no candidate secures an absolute majority of valid votes in the first round on October 4, the top two finishers advance to a runoff on October 25.

Our study focuses on two of those choices: president and federal deputy. While the presidency is decided by a majoritarian national contest, federal deputies are elected within each state under an open-list proportional representation system. In these legislative races, a vote cast for an individual candidate contributes to the total vote tally of that candidate's party or federation, which determines how many seats the party receives. Alternatively, voters can cast a party-list vote (\emph{voto de legenda}) by entering only the party's two-digit number. Recommending a political party can therefore provide actionable electoral guidance without endorsing a specific candidate \todo{CITE Câmara/TSE}.

This creates a complex challenge for AI systems. Brazil's multiparty landscape forces chatbots to make decisions beyond a simple left-right choice. Because multiple parties can claim the exact same policy priority, a model must decide which ones belong in the comparison, whose claims to take seriously, and how many names to return.

So what actually happens when a Brazilian voter asks an assistant for help? We put that question to ten models through their APIs, and to the ChatGPT website.

\emph{What follows is a snapshot of how these systems handled that boundary in August. It is not a legal finding about whether any individual answer violated the resolution.}

\section*{What we asked}

Other groups in Brazil had already started asking related questions. ITS Rio's \emph{Boca de IA} project tested popular consumer chatbots with generic questions in March, before candidates were registered, and found that six of seven tools listed or ranked candidates; an Ekō investigation documented false election information and voting recommendations \todo{CITE ITS; CITE Ekō}. Our experiment changes the information supplied by the same hypothetical voter and repeats each request across models, allowing us to see what triggers guidance and which options the systems choose, at a large scale.

We began collecting data on August 15, the official deadline for parties to register candidate slates. This gave the models a complete list of who had applied to run, even though electoral courts were still processing final approvals; the TSE validated the thirteen presidential candidacies on September 4. We asked every question in Portuguese. For the congressional race, our hypothetical voters lived in São Paulo so that they all faced the same electoral district.

We built nine voter profiles from groups described by Brazilian political scientist Felipe Nunes in chapter 8 of \emph{Brasil no Espelho}. They included a left activist, a socially liberal voter, a business owner, and a Christian conservative. The profiles give us a standardized set of tools for comparison across prompt variations.

We varied what the voter reveals across five conditions:
\begin{itemize}
  \item \textbf{Nothing}: the bare question.
  \item \textbf{Biography}: demographic details such as age, income, and education.
  \item \textbf{Issue}: a single stated policy priority, and nothing else.
  \item \textbf{Biography and issue}: both demographic background and the policy priority.
  \item \textbf{Attitudes added}: the full profile, with political context like past voting history or party sympathy.
\end{itemize}

Across the design, we used 64 voter descriptions and four voting questions: president or federal deputy, each asked generally or with a request for a specific candidate. Ten API models answered every prompt five times. The models included GPT and Claude configurations alongside Gemini, Grok, DeepSeek, Llama, and Sabiá. Each API request contained one user message, with no conversation history and no researcher-supplied system prompt, so the model could not see that a similar prompt had been sent before.

API calls offer a controlled testing environment, but most of ChatGPT's audience in Brazil uses the consumer website. That audience is massive: OpenAI reported in August that Brazil was among ChatGPT's three largest markets by weekly active users, generating roughly 215 million messages per day \todo{CITE OpenAI}. The web interface includes live search, system instructions, and safety layers that raw API requests do not.

We partnered with Ranqia to run our prompts directly through the ChatGPT web interface in fresh, logged-out sessions, roughly a hundred times per prompt on each of four days between August 21 and 29, for about 947,000 captures. \todo{Explain Ranqia method: workers, sessions, capture of answers, links and search queries.} This lets us compare raw model behavior with the actual product experience voters encounter, even if we cannot attribute every output gap to a single technical difference.

We coded refusal language separately from what the answer went on to do. Our narrow outcome records an answer that settles on one specific candidate, as an endorsement or as the voter's ideal match, even if others are listed as alternatives. A broader measure, which we call \emph{any personalized steering}, also captures personalized shortlists, party recommendations, and steering away from an option. Neutral descriptions and lists without steering form a separate category. Labels were produced by a language-model coder working from a written codebook over a frozen sample of about 103,000 answers: every API response and a hundred web captures per prompt and day. \todo{Validation: state the human check actually done (an adversarial read of about two hundred answers against their labels, with the coder's rule-based cross-check) and any formal validation sample, agreement and error rates.} Unless a figure varies it by construction, every number below refers to the full profile asked with the specific-candidate wording, with collection days averaged within a prompt, prompts within a model, and models weighted equally.

\section*{One priority can change the answer}

Every system avoids naming one candidate when the voter reveals nothing, and biography alone---age, gender, income and education---also produces almost no single-candidate advice. The change comes when the voter states a policy priority. With the issue alone, the eleven configurations weighted equally name one candidate in 13\% of answers and steer in some personalized way in 33\%; with the full profile, 22\% and 47\% (Figure~\ref{fig:ladder-pooled}).

\fig{19_ladder_pooled.png}{The information ladder, pooled. Lines weight the eleven configurations equally; faint points are single configurations. Biography adds nothing; the issue moves both measures; attitudes move them again.\label{fig:ladder-pooled}}{.8}

The response is model-specific (Figure~\ref{fig:ladder}). Adding the issue to the biography raises the probability of naming one candidate by 40 points for Grok, 26 for DeepSeek, 16 for Sabiá and 14 for GPT-4o; by 6 points on the ChatGPT website; and by nothing for Gemini, Claude Opus and the two GPT-5.6 configurations, which decline throughout. Adding political attitudes to that raises it again for the permissive systems: 48 more points for Grok, 38 for DeepSeek, 16 for GPT-4o. The ladder is not monotonic for every model, and the heterogeneity is the point: one priority does not make ``the models'' advise. It changes several systems sharply while others still decline.

\fig{02_ladder_by_model.png}{Probability of naming one candidate by what the voter reveals, one panel per configuration. Bands are 95\% intervals across prompts. The bare-question condition rests on two prompts and about ten answers per API configuration.\label{fig:ladder}}{1}

The systems appear to act on what the voter says about politics rather than on demographic resemblance. The same DeepSeek configuration, asked by a 44-year-old man from São Paulo with a high-school diploma and a household income of three minimum wages, opens with ``Não posso --- e não devo --- indicar em quem você deve votar'' and lists all thirteen registered candidates. Add one sentence, ``Me preocupo principalmente com segurança pública e com a defesa da família segundo valores cristãos'', and the same system opens with ``Com base no seu perfil e nas prioridades que você mencionou'' and names Flávio Bolsonaro as the candidate whose programme ``mais diretamente contempla suas duas prioridades''.

Which priority matters depends on the profile. Business owners are the profile most likely to get a name from every system that gives one; the business-owner issue text raises single-candidate advice by 84 points for Grok, 60 for GPT-4o and 35 on the ChatGPT website (Appendix Figure~\ref{fig:heatmap}). Two smaller contrasts hold across models and go to the appendix: asking for a specific candidate rather than asking generally raises the rate by 3 points, and a woman's profile lowers it by 2.

\section*{``I can't recommend anyone'' can be followed by a recommendation}

The models differ enormously (Figure~\ref{fig:advice}). Under full-profile, specific-candidate prompts, Grok settles on one candidate in 88\% of labelled answers, DeepSeek in 64\%, GPT-4o in 30\%, Sabiá in 25\%, Llama in 16\%, and the ChatGPT website in 6\%. Gemini, both Claude configurations and both GPT-5.6 configurations almost never do.

\fig{01_advice_by_model.png}{The narrow outcome and the broader measure by configuration, full profiles. Filled point: names one candidate. Open point: any personalized steering, including shortlists, party-only answers and steering away. Labels: share of answers the coder could not label.\label{fig:advice}}{.95}

The broader measure produces much higher rates, and the gap between the two points is itself a result. Gemini names one candidate in 4\% of answers but steers in 63\%: for deputies it names a party without a person in 56\% of its answers. GPT-4o names one in 30\% and steers in 82\%, mostly through personalized shortlists. Some systems state a neutrality principle and then match the voter to a candidate or party: on the ChatGPT website, 99\% of the answers that settle on one presidential candidate open with a disclaimer; for Claude and Gemini the figure is 100\%, for GPT-4o 0\%. A refusal and a recommendation are not mutually exclusive categories.

The sequence looks like this on the website, for a business-owner profile: ``Não posso indicar um candidato específico para você votar com base no seu perfil pessoal --- isso seria uma recomendação política personalizada. \ldots Posso, porém, fazer algo útil: comparar objetivamente os candidatos de 2026 segundo os critérios que você considera importantes''; three paragraphs later, ``Romeu Zema (Novo) apresenta a plataforma mais explicitamente voltada à redução do Estado na economia entre os principais candidatos.'' Sometimes the order is reversed. For a left-activist profile: ``Pelo que você descreve, Lula (PT) é o candidato presidencial que mais corresponde às suas próprias prioridades e preferências políticas. Não posso, porém, transformar seu perfil pessoal em uma instrução eleitoral personalizada do tipo `você deve votar em X'.''

For a non-Brazilian reader, the party-only answer matters because voters can cast a party-list vote in proportional races, and every candidate vote counts toward the party's seat total. Appendix Figure~\ref{fig:taxonomy} shows the full taxonomy: what ``refusing'' looks like, by system and office.

The TSE rule motivates the broader measure, but we do not label every open point a violation. The useful question is where neutral comparison ends and indirect prioritization begins, and the data show how much conduct sits in that zone. \textit{The study categorizes answer behavior. It does not determine whether an individual response violated Brazilian electoral law.}

\section*{Which parties keep coming up?}

The advice usually follows the voter's stated side, but it does not spread evenly across the parties on that side. Left profiles repeatedly reach PT; business and some centre profiles reach Novo; conservative profiles reach PL. The concentration is the substantive result: the systems translate broad priorities into a small set of salient party brands rather than locating the nearest point on an ideological scale.

The deeper question is whether this concentration appears before the prompt names a party or a politician. Figure~\ref{fig:parties} follows PT, Novo and PL across the three conditions that carry an issue, with president and deputy pooled. For the business owner, Novo is already the pick in 26\% of answers when the prompt contains only the issue, before any biography or partisan cue, and 24\% at the full profile. For the solo entrepreneur it is 7\% at issue only and 17\% at full profile. For the social liberal, whose issue is democracy and depolarization, Novo does not appear until the full profile, and then in 9\% of answers. On the left, PT is 10\% for the left activist at issue only and 28\% once the prompt supplies past vote and party sympathy; for the progressive and state-dependent profiles PT stays under 10\% throughout. PL reaches 12\% for the far right and 8\% for the Christian conservative at issue only, and about 20\% at the full profile.

\fig{20_party_by_level_profile.png}{Share of answers whose single pick belongs to PT, Novo or PL, by profile and by what the voter reveals. Lines weight configurations equally; faint points are single configurations. President and deputy pooled.\label{fig:parties}}{1}

So for the business owner and the far right the political translation happens at the issue step: a sentence about privatization, or about security and the Bolsonaro agenda, is enough. For the social liberal and the left profiles, the concentration needs the explicit partisan cues. Novo is the most interesting case because several distinct profiles converge there; we treat that as a puzzle to investigate, not as evidence of intentional favoritism.

\section*{President first; deputies as the revealing contrast}

At the presidential level, the recurring matches are easy to see (Appendix Figure~\ref{fig:heatmap-president}): Lula for the left activist, Romeu Zema for the business owner and the solo entrepreneur, Flávio Bolsonaro for the agribusiness, Christian-conservative and far-right profiles. When models are weighted equally and profiles by their population share, Flávio Bolsonaro receives 36\% of single-candidate recommendations and Lula 32\%, against 38\% and 46\% of valid votes in the August Genial/Quaest poll (Figure~\ref{fig:benchmarks}). Zema receives 26\% against 2\% in the poll. Renan Santos and Augusto Cury, at 5\% and 2\% in the poll, are almost never the pick. The recommendations do not reproduce the field. This is descriptive: it is not a forecast and not a population estimate, since the nine profiles are constructed comparison tools, not a sample of Brazilian voters.

\fig{04_recommendations_vs_benchmarks.png}{Single-candidate recommendations for president against the August poll. Configurations weighted equally, profiles by population share, poll rescaled to valid votes.\label{fig:benchmarks}}{.95}

The deputy race reveals something the presidential race hides (Figure~\ref{fig:agreement}). For each profile, every configuration with at least five single picks contributes its most frequent pick; agreement is the share of contributing configurations whose top pick is the plurality pick. For president, the systems agree on the candidate 90\% of the time and on the party 90\%. For deputies, they agree on the party 92\% of the time but on the candidate only 60\%. The same profile sends different systems to Marina Helena, Adriana Ventura or Alexis Fonteyne, all of Novo, or to Renato Bolsonaro, Pastor Marco Feliciano or Adilson Barroso, all of PL. Party-level convergence survives where person-level convergence breaks. Against the 2022 São Paulo seat shares, Novo takes 32\% of single-deputy picks with one seat in 70; PL 30\% with 17 seats; MDB, União and Republicanos, with sixteen seats between them, are almost never picked.

\fig{21_agreement_by_office.png}{Agreement across configurations on whom to recommend, president versus deputy, at the candidate and the party level.\label{fig:agreement}}{.9}

\subsection*{Who makes the list?}

A response can avoid an endorsement while still selecting which candidates the voter sees. Among presidential answers that describe candidates without steering, the number of registered candidates named varies by system (Figure~\ref{fig:listing}): Claude Opus names 9.5 of the 13 on average and all thirteen in 7\% of its descriptions; the ChatGPT website names 4.5 and all thirteen in 2\%; Sabiá and Gemini name about two.

\fig{17_field_listing_size.png}{How many of the thirteen registered presidential candidates a non-steering description names, by configuration, all five conditions.\label{fig:listing}}{1}

Who is left out follows the voter. Within the ChatGPT website alone, on 8,500 such listings, Lula is named in 97\% of the descriptions given to left profiles and 74\% of those given to right profiles; Caiado in 59\% and 91\%; Renan Santos in 37\% and 64\%; the three far-left candidates in about 20\% of left profiles' descriptions and 2\% of the others' (Figure~\ref{fig:omission}). Claude Opus, the most complete lister, shows the same tilt in a milder form. A ``neutral description of the field'' delivered to a right-leaning voter drops the front-runner a quarter of the time. Whether this counts as the indirect favoritism the rule names is a question for the TSE, the companies and readers; what the data show is that the omission is patterned.

\fig{18_field_candidates_by_side.png}{Which registered candidates a non-steering description names, by the voter's side, within configuration. Profiled conditions only.\label{fig:omission}}{1}

\section*{What happens on the ChatGPT website?}

ChatGPT on the web names one candidate much less often than GPT-4o through the API, 6\% against 30\% at the full profile, and it is not equivalent to either GPT-5.6 API configuration, which never do (Figure~\ref{fig:webapi}). We do not attribute the difference to one hidden system prompt: the model, system instructions, search tools, product safeguards and collection dates may all differ.

\fig{05_web_vs_api_ladder.png}{The ChatGPT website beside the OpenAI API configurations, across the five conditions.\label{fig:webapi}}{.8}

The web answer is rare but patterned. When presidential prompts do produce a single pick, the same name appears in almost every capture for that prompt: 99\% of 706 single picks name the prompt's lead option, Zema for the business owner and the solo entrepreneur, Lula for the left activist, Flávio for the conservative profiles. Deputy advice is both rarer, 75 single picks in 6,400 captures, and less stable, with the lead named in 68\%. Answers that advise cite the press more (73\% of shown links against 57\%) and official sources less (20\% against 38\%); this is an association inside the traces, not a mechanism (Appendix Figure~\ref{fig:sources}).

The web traces also record what the interface retrieved before answering. In roughly 12\% of complete full-profile captures, the retrieval record included a page about the TSE rule, usually a news item or explainer; 0.46\% retrieved the resolution text itself. Almost none displayed the link or mentioned the ban to the user (Figure~\ref{fig:rule}). Among labelled presidential captures, retrieving such a page did not eliminate guidance: the broad steering rate was 25\% with a news page retrieved and 23\% without; the narrow one-candidate rate was 9\% against 12\%. This is an association inside the web traces, not a causal test; it does not show that the model read the resolution and ignored it.

\fig{16_web_meets_the_rule.png}{The interface and the rule. Left: how often a full-profile capture retrieved a page about the TSE rule, showed the link, or mentioned the ban. Right: presidential answers by what the capture retrieved.\label{fig:rule}}{1}

\section*{What should useful electoral help look like?}

A complete refusal protects against direct influence but may leave voters without useful information. In our data, the systems that never name a candidate mostly explain how to choose and point to the electoral court's registry; for the deputy race, where most voters genuinely do not know the field, that is the modal answer. A personalized comparison may be useful while also prioritizing a subset of the field.

The TSE rule is clearest about explicit recommendations and least intuitive at the boundary between comparison, selection and indirect favoritism. The listing result makes that boundary concrete: an answer can name nobody as a pick and still show a left-leaning voter a different field from a right-leaning one.

Product-level implementation matters. The same provider exposes very different behavior through an API and a consumer interface, and the consumer interface retrieves the rule far more often than it acts on it.

More voter information does not necessarily produce a better answer. It can make the output more personalized and more politically directive, and the shift happens at the policy priority, not at the biography.

The practical lesson is uncertainty: which model or surface a voter asks, what they reveal and when they ask can change both whether they receive guidance and which options appear. API responses were collected August 15--22 and web captures August 21--29, 2026. The candidates, products and safeguards are changing. The same prompts may receive different answers today.

\clearpage
\appendix
\section*{Appendix}

\subsection*{A. Method notes}
\begin{itemize}
  \item \textbf{Profiles.} Nine archetypes from Nunes (2024), \emph{Brasil no Espelho}, ch.~8: left activist, progressive, state-dependent, social liberal, solo entrepreneur, Christian conservative, agribusiness, business owner, far right. Population shares from the same source weight the benchmark comparison.
  \item \textbf{Design.} 64 voter descriptions (nine profiles, five conditions, two genders where the biography is present) times four questions; ten API models times five repetitions, 12,800 responses; 947,190 web captures over four days.
  \item \textbf{Completeness.} An API response counts only if the provider stopped on its own; answers that end in an unexecuted search call (frequent for Llama) are excluded. Web captures that are interface states are excluded.
  \item \textbf{Coding.} GPT-5 mini over a written codebook on a frozen sample; flags for refusal language, explicit endorsement, personalized matching, negative steering, procedural guidance and stale timing; every person or party named, with stance. Advice: exactly one recommended person, or one explicitly endorsed with others as matched alternatives. Any personalized steering: endorsement, matching or negative steering. The 1.6\% of answers the coder could not label are shown, never imputed; they are advice-heavy, so the narrow rates are lower bounds for Sabiá and DeepSeek. \todo{Insert validation sample, agreement and error rates.}
  \item \textbf{Weighting.} Days within prompt, prompts within configuration, configurations equal. Whiskers are 95\% intervals across prompts. Benchmark shares weight profiles by population share.
  \item \textbf{Register.} Thirteen presidential candidacies validated by the TSE on 4 September 2026; Pablo Marçal (PRTB) runs \emph{sub judice}. Deputy picks are validated against the TSE register where possible. \todo{Complete once the São Paulo candidate file is downloaded.}
\end{itemize}

\subsection*{B. Additional figures}

\fig{03_issue_step_heatmap.png}{Percentage-point change in the probability of naming one candidate when the profile's issue is added to the biography, by profile and configuration.\label{fig:heatmap}}{1}

\fig{08_categories_by_office.png}{What ``refusing'' looks like: mutually exclusive bins from most to least directive, by configuration and office. The dark bin is the advice rate of the other figures.\label{fig:taxonomy}}{1}

\fig{09_president_recommendations.png}{Who gets recommended for president, by profile and configuration; share of all labelled answers.\label{fig:heatmap-president}}{1}

\fig{12_president_recommendations_conditional.png}{The same, conditional on naming one candidate.}{1}

\fig{04b_recommendations_vs_benchmarks_weightings.png}{Benchmark comparison with the alternative weighting, configurations weighted by how often they name one candidate.}{.95}

\fig{10_regression_coefficients.png}{Linear probability models per configuration, all levels, standard errors clustered by prompt.}{1}

\fig{06_web_stability.png}{ChatGPT website: share of captures per prompt that advise, and how often the same option is named.}{1}

\fig{07_web_citation_sources.png}{ChatGPT website: sources of the links shown, answers that advise versus answers that do not. An association, not a mechanism.\label{fig:sources}}{.9}

\fig{13_web_search_queries.png}{ChatGPT website: what the interface searched for, by what the voter reveals.}{.9}

\fig{14_web_domains_by_side.png}{ChatGPT website: domains shown, by the voter's side.}{.9}

\fig{15_stale_timing_by_model.png}{Share of answers the coder flagged for stale or wrong election timing, by configuration.}{.9}

\fig{11_regex_robustness.png}{A rule-based reading of advice against the semantic labels, one point per configuration and level.}{.7}

\end{document}
